This deliverable has established a structured catalogue of fourteen use cases for the TenSyGrid Toolbox. UC-01 to UC-05 represent five workflows already prototyped in the Stability Monitor: executing a small-signal stability analysis, screening the system status and eigenvalue spectrum, diagnosing critical modes, checking the power-flow voltage profile, and assessing the small-signal impact of $N-1$ contingencies.

UC-06 to UC-14 define nine proposed extensions. They cover comparative scenario assessment, device-level participation analysis, offline and online tracking of stability margins, critical-element ranking, separate guidance for controller retuning and new-device placement, renewable-integration studies, automated reporting, centre-manifold analysis, and region-of-attraction estimation. The proposals combine the objectives of the TenSyGrid project, experience with existing stability tools, lessons from the current prototype, and evidence gathered through industry interviews.

The next step is to prioritise the proposed use cases within the consortium according to technical feasibility, expected industrial value, and their dependence on the iMTI formulation. The selected cases will then be implemented and validated iteratively with industry stakeholders. Workflows that are not part of the current implementation scope remain documented as possible future extensions rather than as commitments of this deliverable.
